\section{Machine translation and multilingual NLP}
\label{ch:mt}

This is Professor Koehn's field. The goal of this chapter is that you can
speak it correctly: know the data, the evaluation problems, the history of
cross-lingual alignment, and where the ``think in English'' debate stands.

\subsection{A short history, so the vocabulary makes sense}

Statistical machine translation (SMT, 1990s to about 2015) learned
phrase-to-phrase translation probabilities from parallel text and combined
them with a language model; Koehn's 2010 textbook is the reference. Neural MT
(2014 onward) replaced this with a single encoder--decoder network trained
end to end; attention was invented for it (Bahdanau et al.\ 2015), and the
transformer (Vaswani et al.\ 2017) was first a translation model. Koehn's 2020
book covers this era. LLMs are decoder-only transformers trained on raw text
that translate as a side effect of scale; whether they are ``truly
multilingual'' is the question the proposal opens with.

\subsection{Parallel data}

A \emph{parallel corpus} pairs sentences with their translations (Europarl,
UN, OPUS). An \emph{n-way parallel} test set has the same source sentences
in many languages, so any pair of languages is aligned through the source.
Two matter here:

\begin{description}[leftmargin=2.2cm, style=nextline, itemsep=1pt]
\item[NTREX-128] the WMT19 English news test set (1{,}997 sentences, 123 documents) professionally translated into 128 languages. The project's representation grid: first 300 sentences per language for profiles, 1{,}500 for the saved 1,500-sentence study.
\item[FLORES-200] 3{,}001 Wikipedia-derived sentences in 200 languages (dev and devtest). The corpus MEXA reports on; the natural replication set.
\end{description}

Document structure matters: NTREX sentences come in documents, so adjacent
sentences share context. The content-transfer test uses that (the
previous sentence of the same document is the matched context), and any
train/test split must be purged by document, a lesson the misalignment decomposition
taught the project the hard way.

\subsection{Translationese and direction}

Translated text carries the shadow of its source: more literal word order,
source-like lexical choices, simplified style. NTREX is translated
\emph{from} English, so every non-English cell is translationese, and the
proposal itself criticizes benchmarks built by translation. WMT test sets
since 2019 include both directions, which made the project's direct test
possible: measured on native versus translated text for seven languages,
translated text shows a slightly higher language share ($+0.02$ to $+0.06$)
and slightly lower retrieval alignment, while model ordering is preserved at
Spearman 0.99. Say the numbers.

\subsection{Evaluation in MT}

BLEU counts n-gram overlap with a reference; chrF does it over characters
(better for morphologically rich languages); COMET is a learned metric. Human
evaluation is the gold standard and expensive. Three problems carry over to
our benchmarks: translated test sets (translationese and content bias), test
set contamination (text seen in pretraining), and chance floors on
multiple-choice tasks. The project's Belebele protocol is zero-shot
log-likelihood scoring, after discovering that five-shot prompts silently
truncate on models with 2k to 4k contexts.

\subsection{Multilinguality in LLMs}

\begin{description}[leftmargin=3.0cm, style=nextline, itemsep=1pt]
\item[resource tiers] high, mid, low, proxied by web-corpus size (CulturaX, CC-100). Low-resource languages have worse tokenization, worse alignment, and worse task scores, all at once, which is why data size is the confound to remove first.
\item[families and scripts] Indo-European, Sino-Tibetan, Afro-Asiatic, Niger-Congo; Latin, Cyrillic, Devanagari, Arabic, Han. Both predict alignment and performance; both are controls.
\item[curse of multilinguality] at fixed capacity, adding languages dilutes each; deliberately multilingual models (BLOOM, Salamandra, EuroLLM) trade English strength for coverage.
\item[cross-lingual transfer] ability learned in one language showing up in another; parallel data in pretraining is a strong driver (Briakou et al.\ 2023).
\item[tokenizer fairness] Ahia et al.\ (2023), Petrov et al.\ (2023): some languages pay several times more tokens for the same content.
\end{description}

\subsection{Aligning representation spaces across languages}

The classic result (Mikolov et al.\ 2013) is that word-embedding spaces of two
languages are related by a nearly linear map; Conneau et al.\ (2018, MUSE) and
Artetxe et al.\ learned that map without parallel data, and Koehn's group
(Marchisio et al.\ 2020 to 2022) studied when it works. The map is usually
constrained to be orthogonal (a rotation), found by \emph{Procrustes}
analysis. \emph{Hubness} is the high-dimensional nuisance where a few points
are the nearest neighbor of everything; frequent English words are hubs, and
retrieval methods use corrections (CSLS) for it.

The proposal asks whether this still holds inside deep LLM layers. The
project's misalignment decomposition answers with a nested sequence of maps scored on
held-out sentences: cross-language differences are mostly an additive offset
plus a scale (63 to 97 percent), rotation is under one percent, and the
nonlinear remainder is small. The linear-mapping picture survives, in a
degenerate form: mostly a shift.

For sentences rather than words, the standard tools are mean-pooled
representations and retrieval. \textbf{MEXA} (Kargaran et al.\ 2025) asks, for
each non-English sentence, whether its English translation is its nearest
neighbor among $N$ candidates, requiring the true pair to beat both its row
and its column; the score is the fraction of sentences that pass. It depends
on $N$ (harder with more candidates) and uses English as the pivot.
\textbf{Alignment-at-Risk} (AaR@10, this project) takes each sentence's
margin, true-pair cosine minus the best wrong-pair cosine, and averages the
worst decile: a tail measure. \textbf{Information Parity} compares how well a
model compresses a text in language $L$ relative to English. Know how each is
computed so your comparisons are fair.

\subsection{The ``LLMs think in English'' debate}

\begin{description}[leftmargin=3.6cm, style=nextline, itemsep=1pt]
\item[Wendler et al.\ 2024] logit lens on Llama shows English tokens at middle layers for non-English input: ``Llamas work in English''.
\item[Shani and Basirat 2025] middle layers are better described as language-specific plus shared spaces than as internal translation into English.
\item[Tezuka and Inoue 2025] ``transfer neurons'' move representations between shared and language-specific spaces.
\item[Zhao et al.\ 2024; Xu et al.\ 2025] language-specific neurons concentrate in final layers.
\item[Bafna et al.\ 2025 (Murray's group)] intermediate layers are multilingual, dominated by all high-resource languages with English foremost.
\item[Dumas et al.\ 2025] activation patching separates the language of a representation from its concept.
\end{description}

Where LFS sits: it quantifies the language component per layer as a variance
share; the hub test finds a multilingual latent factor beats English as a
predictor for 113 to 124 of 127 languages in every model; and the lens
validity measurement shows why lens-based readouts of the middle are weak.
This supports the Shani/Bafna measure over the ``English'' measure, with
numbers.

\subsection{Improving multilinguality: the Aim 2 options}

The proposal's second aim lists four routes: more parallel and code-switched
data; word- and segment-level alignment losses; an adversarial language
discriminator; multilingual preference training. Prior work: Liu and Niehues
(2025) and Bu et al.\ (2025) align mean-pooled sentence representations in
middle layers and report small gains. The project tested naive forms of the
first three on a 0.6B model: code-switched data gave no detectable benefit
and cost perplexity outside the replay set; the word-level squared-difference
loss collapsed representations by a factor of twenty while retrieval scores
improved; direct LFS minimization deepened the dip through a shortcut. The
lesson for Aim 2: alignment objectives need the raw components as
constraints and independent capability checks. None of this is in the ICLR
paper beyond the stress-test section.

\subsection{Downstream benchmarks you will be asked about}

\begin{description}[leftmargin=2.6cm, style=nextline, itemsep=1pt]
\item[Belebele] 4-choice reading comprehension, 122 languages, translated from FLORES; chance 0.25; 300 questions per language here.
\item[INCLUDE] natively authored benchmark questions in 44 languages (30 overlap the panel); used to test whether the benchmark null was a translation artifact (it was not).
\item[Global-MMLU, MGSM, XNLI] other translated multilingual benchmarks; know they exist and share the translation problem.
\end{description}
